\chapter{زیرساخت توسعه، مدیریت کد، مستندسازی و تحویل‌دادنی‌ها}
\label{chap:devops}
%=====================================================================

\section{ضرورت این فصل}

موفقیت فنی طرح، همان‌طور که در فصل‌های پیشین تشریح شد، به‌تنهایی کافی نیست؛ یک طرح ملی با این سطح از حساسیت امنیتی (زیرساخت بانکی، ارتباط با بانک مرکزی و مرکز افتا) نیازمند یک \textbf{زیرساخت مهندسی نرم‌افزار حاکمیتی} است که چهار بُعد را پوشش دهد: مدیریت سورس‌کد، معماری و زبان‌های برنامه‌نویسی، مستندسازی زنده و رسمی، و کنترل تحویل‌دادنی‌های پروژه. با توجه به تحریم‌های فناوری مؤثر بر دسترسی به سرویس‌های متعارف جهانی (\lr{GitHub}، \lr{Docker Hub}، \lr{GitLab.com})، انتخاب این زیرساخت باید هم‌زمان دو معیار را برآورده کند: \textbf{تداوم دسترسی داخلی} و \textbf{حاکمیت داده و حریم امنیتی} متناسب با ماهیت طرح.

\section{مدیریت سورس‌کد و کنترل نسخه}
\label{sec:vcsgovernance}

با توجه به این‌که کد این طرح شامل پیاده‌سازی الگوریتم‌های رمزنگاری پساکوانتومی و منطق کنترلی \lr{QKD} است، یعنی دقیقاً همان دارایی‌ای که در فصل \ref{chap:security} به‌عنوان هدف حملات احتمالی معرفی شد، میزبانی آن روی سرویس‌های عمومی خارجی از ابتدا منتفی است. دو گزینهٔ داخلی قابل‌بررسی است:

\begin{itemize}
\item \textbf{استقرار \lr{GitLab} سازمانی (\lr{Self-Hosted / On-Premise})}: نصب نسخهٔ \lr{GitLab Enterprise Edition} یا \lr{Community Edition} روی زیرساخت اختصاصی سازمان (یا در دیتاسنتر بانک مرکزی/مجری طرح)، با کنترل کامل بر دسترسی، رمزنگاری در حالت سکون، پشتیبان‌گیری و ممیزی. این گزینه به‌دلیل حاکمیت کامل داده، \textbf{گزینهٔ پیشنهادی برای مخازن هستهٔ امنیتی طرح} (پیاده‌سازی \lr{PQC}، منطق \lr{KMS}، رابط \lr{ETSI GS QKD 014}) است.
\item \textbf{هم‌گیت (\lr{hamgit.ir})}: سرویس رایگان \lr{GitLab} ارائه‌شده توسط هم‌روش که پس از اعمال تحریم‌های \lr{GitLab} برای کاربران ایرانی راه‌اندازی شد و اجراکنندهٔ \lr{CI/CD} با رانرهای داخل کشور را بدون دغدغهٔ تحریم فراهم می‌کند \Cite{hamgit2026}. این گزینه برای فاز اولیهٔ توسعه، مخازن غیرحساس (ابزارهای کمکی، مستندات عمومی، کد نمونه/آموزشی) و به‌عنوان \textbf{پل ارتباطی سریع تا آماده‌شدن زیرساخت سازمانی} پیشنهاد می‌شود.
\end{itemize}

\begin{table}[htbp]
\centering
\caption{مقایسهٔ گزینه‌های میزبانی سورس‌کد}
\label{tab:vcs}
\small
\begin{tabular}{p{3cm}p{5.5cm}p{5.5cm}}
\toprule
\textbf{معیار} & \textbf{\lr{GitLab} سازمانی (روی‌زمین)} & \textbf{هم‌گیت (\lr{hamgit.ir})} \\
\midrule
حاکمیت داده & کامل؛ زیرساخت متعلق به کارفرما & میزبانی نزد شخص ثالث داخلی \\
مناسب برای & مخازن حساس امنیتی/رمزنگاری، \lr{KMS} & توسعهٔ اولیه، ابزارهای جانبی، آموزش \\
\lr{CI/CD} & رانر اختصاصی، قابل ایزوله‌سازی کامل (حتی \lr{Air-Gapped}) & رانر داخل کشور، بدون مشکل تحریم \\
هزینهٔ نگهداشت & نیازمند تیم زیرساخت داخلی & رایگان، نگهداشت برعهدهٔ هم‌روش \\
\bottomrule
\end{tabular}
\end{table}

جریان کاری پیشنهادی، مستقل از گزینهٔ میزبانی: شاخه‌بندی از نوع \lr{Trunk-Based Development} یا \lr{GitLab Flow}، الزام \lr{Merge Request} با حداقل یک بازبین مستقل برای مخازن رمزنگاری، امضای دیجیتال کامیت‌ها (\lr{Signed Commits})، و اجرای خودکار \lr{SAST}\LTRfootnote{Static Application Security Testing}، اسکن مخازن باز (\lr{SCA}) و اسکن نشت اسرار (\lr{Secret Scanning}) در هر \lr{Pipeline}.

\section{بحث کدها: معماری نرم‌افزاری و زبان‌های برنامه‌نویسی}

کدپایهٔ طرح، به‌ویژه در حوزه‌های \lr{QKD} و \lr{PQC}، ماهیتاً چندلایه است و انتخاب زبان باید متناسب با هر لایه صورت گیرد:

\begin{itemize}
\item \textbf{لایهٔ سخت‌افزار/فریمور} (کنترل زمان‌بندی فوتون، رابط آشکارساز تک‌فوتونی، پردازش \lr{FPGA}): \lr{VHDL}/\lr{Verilog} برای منطق زمان‌واقعی سخت‌افزاری، به‌همراه \lr{C} برای درایورهای سطح پایین.
\item \textbf{لایهٔ کنترل و پروتکل \lr{QKD}} (تبادل پایه، تطبیق خطا، تقطیر حریم خصوصی، مدیریت \lr{QBER}): \lr{C++} به‌دلیل نیاز به کارایی بالا، تأخیر کم و کنترل مستقیم حافظه، الزامی برای عملیات زمان‌حساس رمزنگاری.
\item \textbf{لایهٔ ادغام \lr{PQC}}: تکیه بر کتابخانهٔ متن‌باز \lr{liboqs} از پروژهٔ \lr{Open Quantum Safe} \Cite{openquantumsafe} که پیاده‌سازی مرجع \lr{ML-KEM}\Cite{fips203}، \lr{ML-DSA}\Cite{fips204} و \lr{SLH-DSA}\Cite{fips205} را به زبان \lr{C}/\lr{C++} ارائه می‌دهد و از طریق لایهٔ \lr{oqs-provider} با \lr{OpenSSL} ادغام می‌شود؛ این رویکرد از بازتولید ناقص یا نادرست الگوریتم‌های رمزنگاری استاندارد جلوگیری می‌کند.
\item \textbf{لایهٔ \lr{KMS} و رابط \lr{ETSI GS QKD 014}}: پیاده‌سازی سرویس‌های \lr{REST API} با \lr{C++} برای هستهٔ رمزنگاری و \lr{Go} یا \lr{Python} برای لایهٔ میکروسرویس/ارکستراسیون، بسته به الزامات کارایی هر سرویس بانکی.
\item \textbf{لایهٔ شبیه‌سازی، آزمون و تحلیل داده}: \lr{Python} به‌عنوان زبان اصلی، برای شبیه‌سازی پروتکل‌های \lr{BB84}/\lr{E91} پیش از استقرار سخت‌افزار (با ابزارهایی مانند \lr{Qiskit} یا \lr{QuNetSim})، تحلیل آماری \lr{QBER} و نرخ کلید، اتوماسیون آزمون‌های \lr{CI}، و اسکریپت‌های ارکستراسیون استقرار.
\end{itemize}

\begin{table}[htbp]
\centering
\caption{برآورد سهم تقریبی زبان‌های برنامه‌نویسی در کدپایهٔ طرح}
\label{tab:langs}
\small
\begin{tabular}{p{2.5cm}p{2cm}p{9.5cm}}
\toprule
\textbf{زبان} & \textbf{سهم تقریبی} & \textbf{کاربرد اصلی} \\
\midrule
\lr{C++} & $\sim45\%$ & هستهٔ رمزنگاری \lr{PQC}، منطق پروتکل \lr{QKD}، \lr{KMS}، ماژول‌های زمان‌حساس \\
\lr{Python} & $\sim30$ تا $35\%$ & شبیه‌سازی، تحلیل داده و \lr{QBER}، اتوماسیون آزمون، اسکریپت‌های \lr{CI/CD} و ارکستراسیون \\
\lr{VHDL}/\lr{Verilog} & $\sim10\%$ & فریمور \lr{FPGA} کنترل زمان‌بندی فوتون و رابط آشکارساز \\
\lr{Go}/\lr{Rust} & $\sim5$ تا $10\%$ & میکروسرویس‌های پشتیبان، \lr{API Gateway}، ابزارهای زیرساختی امن‌حافظه \\
\bottomrule
\end{tabular}
\end{table}

انتخاب \lr{C++} به‌عنوان زبان غالب هستهٔ رمزنگاری و کنترل، ریشه در نیازهای صنعت \lr{QKD}/\lr{PQC} در سطح جهانی دارد (خود \lr{liboqs} و اغلب پیاده‌سازی‌های مرجع \lr{NIST} به \lr{C}/\lr{C++} نوشته شده‌اند)؛ \lr{Python} در مقابل، به‌دلیل سرعت توسعه و اکوسیستم غنی علمی/داده، برای لایه‌های غیر زمان‌حساس ترجیح داده می‌شود. توصیه می‌شود از \lr{Rust} برای ماژول‌های جدیدی که امنیت حافظه در آن‌ها بحرانی است (مانند تجزیهٔ ورودی‌های نامعتبر در رابط \lr{API}) نیز به‌عنوان گزینهٔ مکمل \lr{C++} استفاده شود.

\section{مدیریت بسته، مخزن آرتیفکت و دستیار هوش مصنوعی توسعه}

برای مدیریت وابستگی‌ها (بسته‌های \lr{C++}/\lr{Conan}، \lr{Python}/\lr{PyPI}، \lr{Go Modules}) و آرتیفکت‌های بیلدشده، دو گزینهٔ استاندارد صنعتی قابل‌استقرار سازمانی وجود دارد:

\begin{itemize}
\item \textbf{\lr{Nexus Repository}}: گزینهٔ سبک‌تر، متن‌باز در نسخهٔ پایه، مناسب برای پروکسی‌کردن و کش‌کردن مخازن عمومی (\lr{PyPI}، \lr{npm}، \lr{Maven}) و میزبانی مخازن داخلی؛ پیشنهاد می‌شود به‌عنوان \textbf{گزینهٔ پیش‌فرض فاز پایلوت} به‌کار رود.
\item \textbf{\lr{JFrog Artifactory}}: گزینهٔ سازمانی‌تر با پشتیبانی گسترده‌تر از انواع بسته، یکپارچگی امنیتی پیشرفته‌تر (\lr{JFrog Xray} برای اسکن آسیب‌پذیری) و مقیاس‌پذیری بالاتر؛ در صورت تخصیص بودجهٔ اشتراک سازمانی، برای مقیاس ملی (فصل \ref{chap:cost}) گزینهٔ مناسب‌تری است.
\end{itemize}

به‌عنوان مکمل، به‌کارگیری \lr{Claude} (دستیار هوش مصنوعی توسعه، از جمله در قالب \lr{Claude Code}) برای تسریع بازبینی کد، تولید مستندات فنی و کمک به تشخیص الگوهای ناامن در کد پیشنهاد می‌شود؛ با این‌حال، با توجه به ماهیت رمزنگاری حساس کدپایه، الزامی است که این ابزار تحت \textbf{قرارداد سازمانی با کنترل حاکمیت داده} (عدم ذخیرهٔ کد برای آموزش مدل، امکان استقرار در محدودهٔ شبکهٔ سازمانی) به‌کار گرفته شود و برای مخازن با بالاترین طبقهٔ محرمانگی (مانند پیاده‌سازی پروتکل مورد اعتماد یا کلیدهای اصلی)، استفاده از دستیارهای هوش مصنوعی متصل به اینترنت عمومی محدود یا ممنوع گردد. سایر ابزارهای پیشنهادی این لایه عبارت‌اند از \lr{SonarQube} برای تحلیل مستمر کیفیت و بدهی فنی کد، و \lr{clang-tidy}/\lr{Coverity} برای تحلیل ایستای تخصصی کد \lr{C++} رمزنگاری.

\section{مخزن کانتینر و میرور تصاویر داکر}

با توجه به محدودیت دسترسی مستقیم به \lr{Docker Hub} ناشی از تحریم، دو اقدام موازی پیشنهاد می‌شود:

\begin{enumerate}
\item \textbf{استقرار مخزن کانتینر سازمانی} (مانند \lr{Harbor}) روی زیرساخت داخلی برای نگهداری تصاویر بیلدشدهٔ اختصاصی طرح، با اسکن آسیب‌پذیری خودکار پیش از استقرار.
\item \textbf{استفاده از میرورهای هم‌داکر (\lr{hamdocker.ir})} به‌عنوان کش‌رجیستری برای دریافت تصاویر پایهٔ عمومی (\lr{hub.hamdocker.ir} به‌عنوان میرور \lr{Docker Hub}، و همچنین میرورهای \lr{mcr}، \lr{gcr}، \lr{quay} و \lr{elastic}) بدون مواجهه با قطعی یا کندی ناشی از تحریم \Cite{hamdocker2026}.
\end{enumerate}

\section{مستندسازی: \lr{Confluence} به‌همراه \lr{LaTeX}}

پیشنهاد می‌شود دو ابزار مکمل، نه رقیب، برای مستندسازی به‌کار روند:

\begin{itemize}
\item \textbf{\lr{Confluence}} برای مستندسازی زنده و روزآمد: تصمیمات معماری (\lr{ADR}\LTRfootnote{Architecture Decision Record})، صورت‌جلسات، راهنمای فنی داخلی، و ویکی دانش تیم؛ با استقرار به‌صورت سازمانی (\lr{Data Center}/داخلی) برای سازگاری با الزامات حاکمیت داده مطرح‌شده در بخش \ref{sec:vcsgovernance}.
\item \textbf{\lr{LaTeX}} برای تحویل‌دادنی‌های رسمی: گزارش‌های فنی نهایی، مقالات علمی، و همین سند پروپوزال، که به همان شیوهٔ این مخزن (ساختار فصل‌بندی‌شده، مدیریت مراجع با \lr{BibTeX}) نگهداری می‌شود. پیشنهاد می‌شود این مخازن \lr{LaTeX} نیز مانند کد نرم‌افزاری تحت کنترل نسخه قرار گیرند و از طریق \lr{CI} (با \lr{latexmk} یا \lr{Tectonic}) به‌صورت خودکار به \lr{PDF} کامپایل و به‌عنوان \lr{Artifact} در دسترس ذی‌نفعان قرار گیرند.
\end{itemize}

\section{بررسی و به‌کارگیری خدمات ابری هم‌روش}

هم‌روش به‌عنوان یکی از ارائه‌دهندگان اصلی زیرساخت ابری داخلی مبتنی بر کوبرنتیز، مجموعه‌ای از خدمات مرتبط با این فصل را ارائه می‌دهد که شامل پلتفرم ابری داکوب (استقرار و مدیریت اپلیکیشن کانتینری)، کوبرنتیز مدیریت‌شده، پایگاه‌دادهٔ مدیریت‌شده، مانیتورینگ و جمع‌آوری لاگ، رجیستری کانتینر اختصاصی، هم‌گیت و هم‌داکر است \Cite{hamravesh2026}. با توجه به این مجموعه، پیشنهاد به‌کارگیری هدفمند به شرح زیر است:

\begin{itemize}
\item \textbf{مناسب برای استفاده روی هم‌روش}: محیط‌های توسعه و آزمون (\lr{Dev}/\lr{Staging})، رانرهای \lr{CI/CD}، ابزارهای شبیه‌سازی غیرحساس، مانیتورینگ و جمع‌آوری لاگ زیرساخت غیرمرتبط با کلید رمزنگاری، و میزبانی هم‌گیت/هم‌داکر برای تسریع توسعه بدون دغدغهٔ تحریم.
\item \textbf{نامناسب برای استفاده روی ابر عمومی (حتی داخلی)}: سامانه‌های تولیدی متصل به \lr{HSM}\LTRfootnote{Hardware Security Module}، پایگاه‌دادهٔ کلیدهای \lr{KMS}، و هرگونه مؤلفه‌ای که مستقیماً با زیرساخت بانک مرکزی یا هستهٔ بانکی تعامل دارد؛ این موارد باید مطابق الزامات نظارتی بانک مرکزی و مرکز افتا، صرفاً روی زیرساخت \textbf{اختصاصی روی‌زمین یا ابر اختصاصی سازمانی} (که خود هم‌روش نیز آن را به‌صورت \lr{On-Premise} ارائه می‌دهد) مستقر شوند.
\end{itemize}

\section{تحویل‌دادنی‌های پروژه}

\begin{table}[htbp]
\centering
\caption{تحویل‌دادنی‌های پیشنهادی به‌تفکیک فاز}
\label{tab:deliverables}
\small
\begin{tabular}{p{4cm}p{10cm}}
\toprule
\textbf{فاز} & \textbf{تحویل‌دادنی‌ها} \\
\midrule
طراحی و مطالعات اولیه & سند معماری فنی (\lr{LaTeX})، سند الزامات امنیتی گرهٔ مورد اعتماد، \lr{ADR}های اولیه در \lr{Confluence} \\
توسعهٔ نرم‌افزار & کدپایهٔ منبع‌باز داخلی در مخازن \lr{Git}، مستندات \lr{API} (رابط \lr{ETSI GS QKD 014})، گزارش پوشش آزمون واحد/یکپارچگی \\
راه‌اندازی پایلوت & گزارش تست میدانی نرخ کلید و \lr{QBER}، گزارش ممیزی امنیتی و آزمون نفوذ مستقل، دفترچهٔ راهنمای بهره‌برداری \\
آموزش و انتقال دانش & بستهٔ آموزشی فنی تیم بهره‌بردار، مستندات نگهداشت در \lr{Confluence} \\
گزارش نهایی طرح & گزارش جامع ملی (\lr{LaTeX}، قابل انتشار برای صندوق نوآوری و معاونت علمی)، ارائهٔ اجرایی، نقشهٔ راه توسعهٔ ملی (فصل \ref{chap:roadmap}) \\
\bottomrule
\end{tabular}
\end{table}

\section{جمع‌بندی زیرساخت پیشنهادی}

\begin{table}[htbp]
\centering
\caption{جمع‌بندی زنجیرهٔ ابزار پیشنهادی توسعه و تحویل}
\label{tab:toolchain}
\small
\begin{tabular}{p{4cm}p{10cm}}
\toprule
\textbf{حوزه} & \textbf{ابزار پیشنهادی} \\
\midrule
سورس‌کد و کنترل نسخه & \lr{GitLab} سازمانی (مخازن حساس) + هم‌گیت (\lr{hamgit.ir}) برای فاز اولیه \\
زبان‌های اصلی & \lr{C++} (هستهٔ رمزنگاری/کنترل)، \lr{Python} (شبیه‌سازی و ارکستراسیون)، \lr{VHDL}/\lr{Verilog} (فریمور) \\
کتابخانهٔ \lr{PQC} & \lr{liboqs}/\lr{oqs-provider} (\lr{Open Quantum Safe}) \\
مدیریت بسته/آرتیفکت & \lr{Nexus Repository} (پایلوت) یا \lr{JFrog Artifactory} (مقیاس ملی) \\
دستیار هوش مصنوعی توسعه & \lr{Claude} با قرارداد سازمانی و کنترل حاکمیت داده \\
مخزن و میرور کانتینر & \lr{Harbor} سازمانی + میرورهای هم‌داکر (\lr{hamdocker.ir}) \\
مستندسازی & \lr{Confluence} (زنده) + \lr{LaTeX} (رسمی/تحویل‌دادنی) \\
زیرساخت ابری غیرحساس & خدمات هم‌روش (داکوب، کوبرنتیز مدیریت‌شده، مانیتورینگ) \\
زیرساخت حساس/تولیدی & روی‌زمین یا ابر اختصاصی سازمانی، مطابق الزامات بانک مرکزی و مرکز افتا \\
\bottomrule
\end{tabular}
\end{table}

%=====================================================================
